\section*{Analysis of features 3 and 4}

\subsection*{What do you observe?}
Observing the scatter plot (\texttt{features\_3\_4\_scatter.png}), we can analyze the distribution of the two classes ("Fake" and "Genuine"). [Insert visual observation here: e.g., The classes appear to occupy distinct regions, or they are heavily mixed. Note which class tends to have higher/lower values for Feature 3 and Feature 4].

\subsection*{Do the classes overlap? If so, where?}
Yes, the classes exhibit overlap. According to the computed data in \texttt{lab02\_project\_feature\_overlap\_ranges.csv}:
\begin{itemize}
    \item \textbf{Feature 3} overlaps between [Insert Min Overlap] and [Insert Max Overlap] (an overlap width of [Insert Overlap Width]).
    \item \textbf{Feature 4} overlaps between [Insert Min Overlap] and [Insert Max Overlap] (an overlap width of [Insert Overlap Width]).
\end{itemize}

\subsection*{Do the classes show similar mean for these two features?}
[Insert Yes/No based on CSV data], the means are [distinctly different / relatively similar]. Extracted from \texttt{lab02\_project\_statistics.csv}:
\begin{itemize}
    \item \textbf{Feature 3:} The "Fake" mean is [Insert Fake Mean], while the "Genuine" mean is [Insert Genuine Mean].
    \item \textbf{Feature 4:} The "Fake" mean is [Insert Fake Mean], while the "Genuine" mean is [Insert Genuine Mean].
\end{itemize}

\subsection*{Are the variances similar for the two classes?}
The variances [show a significant difference / are relatively similar]. According to \texttt{lab02\_project\_statistics.csv}:
\begin{itemize}
    \item \textbf{Feature 3:} The variance for "Fake" is [Insert Fake Variance], and for "Genuine" it is [Insert Genuine Variance].
    \item \textbf{Feature 4:} The variance for "Fake" is [Insert Fake Variance], and for "Genuine" it is [Insert Genuine Variance].
\end{itemize}

\subsection*{How many modes are evident from the histograms?}
Based on \texttt{feature\_3\_histogram.png} and \texttt{feature\_4\_histogram.png}, we can observe the following modes:
\begin{itemize}
    \item \textbf{Feature 3:} The "Fake" distribution exhibits [Insert Number, e.g., one/two] mode(s), and the "Genuine" distribution exhibits [Insert Number] mode(s).
    \item \textbf{Feature 4:} The "Fake" distribution exhibits [Insert Number] mode(s), and the "Genuine" distribution exhibits [Insert Number] mode(s).
\end{itemize}
